% !TeX program = lualatex
% Build with LuaLaTeX twice (TeX Live 2025 or newer), or use latexmk.
%
% ATS / machine-readability notes:
% - Single-column layout; no tables, multicolumns, or text boxes for content.
% - Every link is also printed as plain text, so parsers that ignore
%   hyperlinks still capture it.
% - Standard section headings (Education, Experience, Publications, Skills).
% - Dates use MM/YYYY where the month is known.
% - Tagged PDF/UA output with a real text layer (OpenType font, Unicode).
\DocumentMetadata{
  lang=en-US,
  pdfversion=2.0,
  testphase={phase-III,title,table},
  pdfstandard=ua-2
}
\documentclass[10pt,letterpaper]{article}

\usepackage[
  top=0.6in,
  bottom=0.6in,
  left=0.7in,
  right=0.7in
]{geometry}
\usepackage{libertinus}
\usepackage{microtype}
\usepackage{enumitem}
\usepackage{titlesec}
\usepackage{xcolor}

\definecolor{ink}{HTML}{202428}
\definecolor{muted}{HTML}{5A6168}
\definecolor{rule}{HTML}{B8BEC4}
\definecolor{link}{HTML}{294E6B}

\usepackage[
  colorlinks=true,
  urlcolor=link,
  linkcolor=link,
  pdfauthor={Chris Lee},
  pdftitle={Chris Lee - Curriculum Vitae},
  pdfsubject={Curriculum vitae: Ph.D. candidate in Operations Research},
  pdfkeywords={Chris Lee, operations research, online learning, bandit
    algorithms, sequential decision-making, optimization, Bayesian modeling,
    machine learning, data science, data engineering, Python, SQL, BigQuery,
    Airflow, Google Cloud Platform, applied science, research science,
    internship}
]{hyperref}

\color{ink}
\pagestyle{plain}
\setlength{\parindent}{0pt}
\setlength{\parskip}{0pt}
\setlength{\emergencystretch}{2em}
\urlstyle{same}

\titleformat{\section}
  {\large\bfseries\color{ink}}
  {}
  {0pt}
  {}
  [\vspace{0.15em}{\color{rule}\titlerule}]
\titlespacing*{\section}{0pt}{0.75em}{0.4em}

\setlist[itemize]{
  leftmargin=1.25em,
  itemsep=0pt,
  topsep=0.05em,
  parsep=0pt,
  partopsep=0pt
}
\setlist[enumerate]{
  leftmargin=1.55em,
  labelsep=0.45em,
  itemsep=0.15em,
  topsep=0.1em,
  parsep=0pt,
  partopsep=0pt
}

% One entry = plain lines of text (no tables), so parsers read them in order.
% #1 organization or title, #2 dates, #3 role / degree, #4 location
\newcommand{\entry}[4]{%
  \textbf{#1}\hfill\textbf{#2}\par
  \emph{#3}\hfill #4\par
}

\newcommand{\paperstatus}[1]{{\color{muted}\emph{#1}}}

\begin{document}

% ---------- Header (plain text, no table) ----------
{\Huge\bfseries Chris Lee}\par
\vspace{0.3em}
{\large Ph.D. Candidate in Operations Research, University of Michigan}\par
\vspace{0.3em}
Ann Arbor, MI \enspace|\enspace (206) 850-9429 \enspace|\enspace
\href{mailto:chrisor@umich.edu}{chrisor@umich.edu}\par
Website: \url{https://leechr001.github.io}
\enspace|\enspace
GitHub: \url{https://github.com/leechr001}\par
Google Scholar:
\url{https://scholar.google.com/citations?user=eetYVLoAAAAJ}\par
\vspace{0.2em}
\textbf{Expected graduation:} 05/2029
\enspace|\enspace
\textbf{Available:} Summer 2027 (05/2027 -- 08/2027)\par

\section*{Research Interests}

Online learning and sequential decision-making under structured feedback;
bandit algorithms; Bayesian models of sequential decisions; applications in
workforce management, energy, and healthcare.

\section*{Education}

\entry
  {University of Michigan}
  {09/2024 -- 05/2028 (expected)}
  {Ph.D. in Operations Research, Department of Industrial and Operations Engineering}
  {Ann Arbor, MI}
Advisors: Xiuli Chao and Izak Duenyas. Research area: online learning and
stochastic bandits.\par

\vspace{0.45em}
\entry
  {University of Washington}
  {2018 -- 2022}
  {B.S. in Mathematics with Honors}
  {Seattle, WA}
GPA: 3.7. Distinctions: Departmental Honors; Dean's List (2018--2022).\par
Honors thesis: \emph{Computing Mixing Times of Parallel Tempered Markov
Chains}. Advisors: Soumik Pal and Rekha Thomas.\par

\section*{Industry Experience}

\entry
  {AmFam Group (American Family Insurance)}
  {08/2022 -- 08/2024}
  {Cloud Data Engineer}
  {Madison, WI (remote)}
\begin{itemize}
  \item Lead architect and developer, on a five-person team, of a 12-month
    initiative to automate key parts of the actuarial ratemaking process over
    terabyte-scale data (Google Cloud Storage, Cloud Functions, BigQuery,
    Cloud Composer, Airflow, Google Workspace).
  \item Built a business-intelligence semantic layer in BigQuery in
    collaboration with data governance engineers, data scientists, data
    analysts, and project managers (Cloud Composer, Airflow, Python, SQL).
  \item Developed an internal tool for evaluating data quality of
    terabyte-scale tables.
  \item Received the 2023 ``Boomtown Developer'' award from the AmFam IT
    organization for outstanding achievement and passion for learning.
  \item Rated ``exceeds expectations'' after three months and
    ``outstanding'' after one year in performance reviews.
\end{itemize}

\vspace{0.3em}
\entry
  {American Family Insurance}
  {06/2021 -- 08/2021}
  {Data Analytics Intern}
  {Madison, WI (remote)}
\begin{itemize}
  \item Produced monthly KPI reports and ad hoc analyses for the performance
    management department of AmFam Personal Auto lines: pulled data from GCP
    with SQL, analyzed in Excel and R, and presented findings in PowerPoint.
  \item Trained a regression model to estimate key partners' mix of business
    (R).
  \item Researched and pitched a nonprofit to AmFam's Philanthropy team,
    resulting in a \$500 grant.
  \item Rated as ``exceeded expectations'' and a ``fast learner'' in
    performance review.
\end{itemize}

\vspace{0.3em}
\entry
  {Market Games}
  {09/2020 -- 12/2020}
  {Marketing Analytics and Data Science Intern}
  {Berkeley, CA}
\begin{itemize}
  \item Independently built and managed a social media marketing data
    pipeline: batch ingestion of relational and non-relational sources,
    including a web scraper for competitive analysis (Python), joined with
    Google Analytics data.
  \item Delivered weekly reports with actionable insights for a non-technical
    team, using clustering for customer segmentation, sentiment analysis of
    customer interactions, and regression models to optimize marketing
    variables (R).
  \item Increased Instagram followers by 30\% and likes per post by about
    90\%.
\end{itemize}

\section*{Research Experience}

\entry
  {Washington Experimental Mathematics Lab, University of Washington}
  {09/2021 -- 12/2021; 03/2022 -- 06/2022}
  {Undergraduate Researcher, Department of Mathematics}
  {Seattle, WA}
\begin{itemize}
  \item Conducted research on graph boundary structure, Markov-chain mixing,
    graphical designs, and spectral graph theory through an honors thesis,
    the Washington Experimental Mathematics Lab, and the Washington Directed
    Reading Program.
  \item Co-authored a paper published in the \emph{Journal of Combinatorics};
    research acknowledged in a paper on graphical designs and Gale duality.
  \item Faculty mentors: Stefan Steinerberger, Rekha Thomas, and Soumik Pal;
    graduate mentors: Kevin Liu and Catherine Babecki.
\end{itemize}

\vspace{0.3em}
\entry
  {CocoBot, UW School of Public Health}
  {06/2019 -- 12/2019}
  {Research Assistant}
  {Seattle, WA}
\begin{itemize}
  \item Prototyped a conversational healthcare AI with Rasa NLU as part of a
    funded research project (Python).
  \item Scraped websites for medical diagnostic criteria and treatments and
    ingested them into Elasticsearch.
\end{itemize}

\section*{Publications and Working Papers}

{\small\bfseries Working Papers}
\begin{enumerate}
  \item \textbf{Chris Lee}, Xiuli Chao, and Izak Duenyas.
    ``Sequential Hiring of Contingent Workers Through Learning-Based
    Optimization.'' \paperstatus{Submitted to Operations Research}, 2026.

  \item \textbf{Chris Lee}, Xiuli Chao, and Izak Duenyas.
    ``The Price of Shared Shocks: From Coverage to Transport in Combinatorial
    Semi-Bandits.'' \paperstatus{Working paper}, 2026.
\end{enumerate}

\vspace{0.12em}
{\small\bfseries Refereed Publication}
\begin{enumerate}
  \item Nick Chiem, William Dudarov, \textbf{Chris Lee}, Sean Lee, and Kevin Liu.
    ``A Characterization of Graphs with at Most Four Boundary Vertices.''
    \emph{Journal of Combinatorics}, 15(3), 2024.
    arXiv: \href{https://arxiv.org/abs/2209.04438}{2209.04438}.
    \paperstatus{Authors listed alphabetically.}
\end{enumerate}

\vspace{0.12em}
{\small\bfseries Acknowledged In}
\begin{itemize}
  \item ``Graphical Designs and Gale Duality.''
    arXiv: \href{https://arxiv.org/abs/2204.01873}{2204.01873}.
\end{itemize}

\section*{Teaching Experience}

\entry
  {Project Management}
  {2025, 2026}
  {Teaching Assistant, University of Michigan}
  {Ann Arbor, MI}
% Optional: add a bullet on what you did, e.g. leading sections, grading
% team projects, or advising student teams -- only what is accurate.

\vspace{0.3em}
\entry
  {Linear Programming (Graduate Level)}
  {2026}
  {Teaching Assistant, University of Michigan}
  {Ann Arbor, MI}

\vspace{0.3em}
\entry
  {Advanced Linear Algebra Tools and Applications}
  {01/2022 -- 03/2022}
  {Undergraduate Teaching Assistant, University of Washington}
  {Seattle, WA}

\section*{Selected Projects}

\begin{itemize}
  \item Backend for an airline flight search and ticketing application (Java, SQL).
  \item Microtargeting system for political advertisements (R, HTML).
  \item Breaking Google CAPTCHA with machine vision (Python).
  \item AWS Athena to BigQuery SQL conversion tool (SQL, Python, regex).
  \item Multi-platform data profiling (Python, SQL, GCP, Excel).
  \item Home automation with smart appliances (Python).
\end{itemize}

\section*{Skills}

\begin{itemize}
  \item \textbf{Methods:} online learning, multi-armed and combinatorial
    semi-bandits, sequential decision-making, Bayesian modeling,
    optimization, linear programming, Markov chains and Monte Carlo methods,
    spectral graph theory, regression, clustering, sentiment analysis.
  \item \textbf{Programming:} Python (pandas, scikit-learn, PySpark,
    TensorFlow), R (dplyr, tidyverse, ggplot2), SQL (BigQuery, Athena), Java,
    regular expressions.
  \item \textbf{Data and cloud:} Google Cloud Platform, Amazon Web Services,
    Apache Airflow, Apache Spark, Docker, Git, GitLab; ETL, data warehousing,
    data lakes, database design and normalization, star-schema data
    modeling, data governance; NoSQL (DynamoDB, Elasticsearch, Neo4j).
  \item \textbf{Project management:} business requirements documents,
    source-to-target mapping, scoping, release cadence, Agile development
    (Jira, Confluence); cross-functional collaboration with engineers, data
    scientists, analysts, and project managers.
  \item \textbf{Tools:} Excel (pivot tables, XLOOKUP), PowerPoint.
\end{itemize}

\section*{Awards}

\begin{itemize}
  \item ``Boomtown Developer,'' American Family Insurance IT, 2023.
  \item 1st Place, UW CSE 416 Kaggle Competition, 2020 (Python).
  \item 1st Place, Best Data Analysis, UW Datathon, 2019 (R, Excel, PowerPoint).
  \item 1st Place, UW CSE Bot Tournament, 2018 (Java).
  \item Data Science Certificate, DubsTech.
\end{itemize}

% Maintenance notes (not printed):
% - Links for the two working papers are omitted because the website currently
%   contains placeholder href="#" values.
% - The graph paper uses its journal publication year (2024), not the preprint
%   year shown on the website.
% - Expected graduation (05/2028), Ph.D. start (09/2024), and availability
%   window are assumptions -- verify before sending.
% - Base CV for industry applied/data/research science roles. Update the
%   "Available" line each recruiting cycle; tailor per posting only if a
%   posting asks for something specific.
% - Resume vs. CV conflicts to resolve: AmFam title ("Cloud Data Engineer" on
%   resume, "Lead Data Analyst" on old CV); UW TA dates (01-03/2022 on resume,
%   2021 on old CV); honors thesis title (resume lists "A Markov Chain Monte
%   Carlo Algorithm for Finding Graphical Designs").
% - Omitted from resume: LinkedIn skill-assessment certificates, MS Teams,
%   Majesco PAS, Guidewire (low signal for research/science roles).

\end{document}